\documentclass[border=4pt]{standalone}
% Shared extension of the installed skill's publication-grade preamble.
% TraceBank figure setup: packages, colour values and plot defaults.
\usepackage{tikz,pgfplots,amsmath,siunitx}
\pgfplotsset{compat=1.18}
\usetikzlibrary{positioning,calc,arrows.meta,fit,backgrounds,shapes.geometric}
\definecolor{cbBlue}{RGB}{0,114,178}
\definecolor{cbOrange}{RGB}{230,159,0}
\definecolor{cbGreen}{RGB}{0,158,115}
\definecolor{cbVermillion}{RGB}{213,94,0}
\definecolor{cbPurple}{RGB}{204,121,167}
\pgfplotsset{every axis/.append style={font=\small,line width=.6pt,
 tick style={black},grid=major,grid style={gray!25,line width=.3pt},
 legend style={font=\footnotesize,draw=gray!50,fill=white},
 label style={font=\small}}}

\usepackage{fontspec}
\usepackage{xeCJK}
\setmainfont{Arial}
\setCJKmainfont{Microsoft YaHei}
\setsansfont{Arial}
\setCJKsansfont{Microsoft YaHei}
\definecolor{ink}{HTML}{17343E}
\definecolor{muted}{HTML}{6C7C80}
\tikzset{
  every node/.style={font=\fontsize{8.3}{12}\selectfont,text=ink},
  box/.style={draw=cbBlue!45,fill=cbBlue!5,rounded corners=2mm,
    minimum height=11mm,text width=31mm,align=center,inner sep=3mm,line width=.6pt},
  safe/.style={box,draw=cbGreen!65!black,fill=cbGreen!7},
  caution/.style={box,draw=cbOrange!80!black,fill=cbOrange!9},
  strong/.style={box,draw=cbVermillion!80!black,fill=cbVermillion!7},
  neutral/.style={box,draw=muted!45,fill=muted!4},
  arrow/.style={-{Stealth[length=2mm,width=1.5mm]},line width=.65pt,draw=ink!75},
  response/.style={arrow,dashed,draw=cbBlue!70},
  note/.style={align=left,font=\fontsize{7.5}{11}\selectfont,text=muted},
  tinylabel/.style={font=\fontsize{7.2}{9}\selectfont,text=muted,fill=white,inner sep=1mm},
  title/.style={anchor=west,font=\fontsize{13}{17}\selectfont\bfseries,text=ink},
  subtitle/.style={anchor=west,font=\fontsize{8}{11}\selectfont,text=muted},
  boundary/.style={draw=cbBlue!45,dashed,rounded corners=3mm,line width=.65pt}
}

% ===== TUNABLES =====
\def\figwidth{183mm}
\def\diagramtitle{F3 · 红色转账的授权与资金时序}
% Source: API /chat, /operations/{id}/authorize, /auth/demo-inbox; BankService
\begin{document}
\begin{tikzpicture}[x=10mm,y=10mm]
\useasboundingbox (0,-14.6) rectangle (18.3,1.5);
\node[title] at (0,1.2) {\diagramtitle};
\node[subtitle] at (0,.6) {示例：向明确联系人转账 1200 元；验证码在模型调用链之外提交};
\foreach \x/\name/\id in {1.5/用户与客户界面/user,5.25/规划与编排器/agent,9.0/API 与 BankService/bank,12.75/独立验证面板/panel,16.5/事务账本/ledger}{
  \node[box,text width=27mm,minimum height=9mm] (\id) at (\x,-.35) {\name};
  \draw[draw=muted!40,dashed] (\x,-.9) -- (\x,-12.8);
}
\draw[arrow] (1.5,-1.25) -- node[tinylabel,above] {自然语言需求} (5.25,-1.25);
\draw[arrow] (5.25,-2.05) -- node[tinylabel,above] {创建待授权操作} (9,-2.05);
\draw[response] (9,-2.85) -- node[tinylabel,above] {红色 / 待确认 / 操作 ID} (5.25,-2.85);
\draw[response] (5.25,-3.65) -- node[tinylabel,above] {展示收款人与金额} (1.5,-3.65);
\draw[arrow] (1.5,-4.45) -- node[tinylabel,above] {明确确认 + 登录会话 + CSRF} (9,-4.45);
\draw[arrow] (9,-5.25) -- node[tinylabel,above] {创建操作绑定挑战} (12.75,-5.25);
\draw[arrow] (12.75,-6.05) -- node[tinylabel,above] {读取模拟收件箱} (9,-6.05);
\draw[response] (9,-6.85) -- node[tinylabel,above] {模拟 OTP，仅在独立面板显示} (12.75,-6.85);
\draw[arrow] (12.75,-7.65) -- node[tinylabel,above] {OTP + 原操作 ID} (9,-7.65);
\node[strong,text width=39mm,minimum height=9mm] at (9,-8.65) {核验 OTP、归属、有效期\\事务内重算风险与幂等状态};
\draw[arrow] (9,-9.55) -- node[tinylabel,above] {一次性原子扣款与入账} (16.5,-9.55);
\draw[response] (16.5,-10.35) -- node[tinylabel,above] {唯一业务回执} (9,-10.35);
\draw[response] (9,-11.15) -- node[tinylabel,above] {更新持久化任务状态} (5.25,-11.15);
\draw[response] (5.25,-11.95) -- node[tinylabel,above] {展示真实回执与完成状态} (1.5,-11.95);
\node[caution,text width=170mm,minimum height=10mm] at (9.1,-13.5) {故障分支：提交后回执丢失 $\rightarrow$ 状态 unknown $\rightarrow$ 查询原操作并对账（reconcile）\\不新建扣款、不把超时当失败重转；重复请求使用原幂等键与操作 ID。};
\node[note,anchor=west] at (.1,-14.4) {认证与业务请求直接进入可信 API；模型不读取 OTP，也不能自行声明强验证已经完成。};
\end{tikzpicture}
\end{document}
